\documentclass{article}

\usepackage[margin=1in]{geometry}
\usepackage{amsmath, amssymb}
\usepackage{graphicx}
\usepackage{booktabs}
\usepackage{multirow}
\usepackage{array}
\usepackage{hyperref}
\usepackage{natbib}
\usepackage{caption}
\usepackage{subcaption}
\usepackage{float}
\usepackage{xcolor}

\hypersetup{colorlinks=true, linkcolor=blue, citecolor=blue, urlcolor=blue}

\title{Can LLMs Identify Cipher Types Zero-Shot?}

\author{Eric Lee \\
  Yee Collins Research Group \\
  \small\texttt{e.lee@ycrg-labs.org}
}

\date{}

\begin{document}
\maketitle

% ============================================================
\begin{abstract}
We evaluate whether large language models can identify classical cipher types
from ciphertext alone, without examples or fine-tuning.
Four models---Claude Sonnet 4.6, GPT-4o, DeepSeek-R1, and Llama 3.3 70B---are
benchmarked across 1,372 samples spanning 9 cipher types and three prompt strategies.
All models exceed the 11.1\% random baseline, with Claude Sonnet reaching 91.3\%
under a multiple-choice prompt.
We find that chain-of-thought prompting \emph{consistently hurts} performance
(collapsing GPT-4o from 74.9\% to 24.3\%), that prompt effects are strongly
model-dependent, and that ROT13 is identified via memorization rather than
reasoning. Nearly all pairwise comparisons are significant after Holm--Bonferroni correction.
\end{abstract}

% ============================================================
\section{Introduction}

Classical cryptanalysis requires recognizing an unknown cipher from ciphertext
alone---a task that has traditionally demanded statistical expertise and domain
knowledge~\citep{sinkov1966elementary}.
Large language models may have absorbed cipher knowledge during pretraining,
raising the question: \emph{can LLMs perform cipher identification zero-shot?}

This provides a clean testbed for studying prompt sensitivity, emergent knowledge,
and the limits of LLM pattern recognition, with practical relevance to automated
security analysis.
We contribute (1) a 1,372-sample benchmark spanning 9 cipher types and 14
configurations, (2) a zero-shot evaluation of four models across three prompt
strategies, and (3) statistical analysis revealing systematic misclassification
patterns and a ROT13 memorization effect.

% ============================================================
\section{Setup}

\paragraph{Dataset.}
We encrypt 98 English plaintexts (43--80 characters) under 14 cipher
configurations across 9 types: Caesar (shifts 3, 7, 19), Atbash,
Vigen\`{e}re (keys \texttt{ACL}/\texttt{CIPHER}/\texttt{SECRET}), ROT13,
Rail Fence (2 and 3 rails), Reverse, Base64, Morse, and Bacon.
Variants of the same cipher share a label, yielding a 9-class problem.
Caesar and Vigen\`{e}re each contribute 294 samples (three variants $\times$ 98),
Rail Fence contributes 196 (two variants), and the remaining six ciphers contribute
98 samples each, for $14 \times 98 = 1{,}372$ total.
Two planned plaintext entries were not collected, accounting for the shortfall from 1,400.
The uniform random baseline is 11.1\%; the majority-class baseline
(always predicting Caesar or Vigen\`{e}re) is 21.4\%.

\paragraph{Models and prompts.}
All models are queried at \texttt{temperature=0}.
We test three prompt strategies: \textbf{Baseline} (``Answer with just the cipher name''),
\textbf{Choices} (same, with the 9 options listed), and \textbf{CoT}
(``Think step by step\ldots Analysis:'').
Responses are parsed with a rule-based extractor handling common aliases;
unparseable responses are scored as incorrect.
Pairwise comparisons use McNemar's test on matched samples; ROT13 vs.\ Caesar
comparisons use chi-squared. All pairwise McNemar tests are corrected jointly
using Holm--Bonferroni. All CIs are 95\% bootstrap over 10,000 iterations.

% ============================================================
\section{Results}

Table~\ref{tab:main} reports overall and per-category accuracy.
All models substantially exceed both baselines.
Claude Sonnet achieves the best overall accuracy of 91.3\% (89.8--92.7\%)
under the choices prompt, with a macro-average accuracy of 94.9\%,
confirming the result is not an artifact of class imbalance.
All pairwise model differences are statistically significant after
Holm--Bonferroni correction ($p_{\text{holm}} < 0.001$, McNemar's test);
full contingency counts, odds ratios, and corrected $p$-values are reported
in Appendix~\ref{app:stats}.

\begin{table}[H]
\centering
\caption{Overall and per-category accuracy (\%) with 95\% bootstrap CIs.
Best result per model in \textbf{bold}.}
\label{tab:main}
\small
\begin{tabular}{llccccc}
\toprule
\textbf{Model} & \textbf{Prompt} & \textbf{Overall} & \textbf{Subst.} & \textbf{Transp.} & \textbf{Encoding} \\
\midrule
\multirow{3}{*}{Claude Sonnet 4.6}
  & Baseline & 87.0 {\scriptsize(85.2--88.8)} & 92.1 & 93.2 & 67.3 \\
  & Choices  & \textbf{91.3 {\scriptsize(89.8--92.7)}} & 85.5 & 98.3 & 99.7 \\
  & CoT      & 73.8 {\scriptsize(71.5--76.1)} & 73.5 & 88.4 & 60.2 \\
\midrule
\multirow{3}{*}{GPT-4o}
  & Baseline & 74.4 {\scriptsize(72.1--76.7)} & 80.1 & 34.0 & 99.7 \\
  & Choices  & \textbf{74.9 {\scriptsize(72.6--77.3)}} & 81.1 & 33.3 & 100.0 \\
  & CoT      & 24.3 {\scriptsize(22.0--26.5)} & 10.1 & 14.6 & 71.8 \\
\midrule
\multirow{3}{*}{DeepSeek-R1}
  & Baseline & \textbf{65.1 {\scriptsize(62.5--67.6)}} & 48.2 & 76.2 & 99.0 \\
  & Choices  & 46.1 {\scriptsize(43.4--48.7)} & 35.6 & 45.9 & 74.1 \\
  & CoT      & 62.8 {\scriptsize(60.2--65.3)} & 47.2 & 74.1 & 92.9 \\
\midrule
\multirow{3}{*}{Llama 3.3 70B}
  & Baseline & \textbf{57.1 {\scriptsize(54.5--59.8)}} & 62.0 & 34.7 & 66.7 \\
  & Choices  & 51.5 {\scriptsize(48.9--54.1)} & 30.7 & 58.5 & 100.0 \\
  & CoT      & 44.9 {\scriptsize(42.2--47.5)} & 38.3 & 46.9 & 60.5 \\
\midrule
\multicolumn{2}{l}{\textit{Random baseline}} & 11.1 & --- & --- & --- \\
\bottomrule
\end{tabular}
\end{table}

\paragraph{Per-cipher breakdown.}
Figure~\ref{fig:heatmap} shows per-cipher accuracy at each model's best prompt.
Encoding ciphers (Base64, Morse, Bacon) are identified near-perfectly by three of
four models, driven by distinctive character sets rather than cryptanalytic reasoning.
Caesar and Rail Fence are the hardest ciphers: GPT-4o scores 0.0\% on Rail Fence
under choices, and Llama scores only 9.9\% on Caesar, systematically defaulting
to Vigen\`{e}re. DeepSeek struggles with Vigen\`{e}re itself (10.2\%), a striking
reversal from Llama's pattern.

\begin{figure}[H]
  \centering
  \includegraphics[width=\linewidth]{fig1_heatmap.png}
  \caption{Per-cipher accuracy (\%) at each model's best prompt. Red = 0\%, green = 100\%.}
  \label{fig:heatmap}
\end{figure}

\paragraph{Confusion patterns.}
Figure~\ref{fig:confusion} shows row-normalized confusion matrices for all models.
Claude Sonnet is nearly diagonal; its main error is Caesar $\to$ ROT13 (25\%),
reasonable since ROT13 is Caesar-13.
GPT-4o completely fails on Rail Fence despite strong overall performance.
Llama exhibits a strong Vigen\`{e}re bias, misclassifying Caesar (64\%),
Rail Fence (88\%), and Bacon (50\%) as Vigen\`{e}re.

\begin{figure}[H]
  \centering
  \begin{subfigure}[t]{0.48\linewidth}
    \includegraphics[width=\linewidth]{fig2_confusion_claude-sonnet.png}
    \caption{Claude Sonnet 4.6 (choices)}
  \end{subfigure}
  \hfill
  \begin{subfigure}[t]{0.48\linewidth}
    \includegraphics[width=\linewidth]{fig2_confusion_gpt-4o.png}
    \caption{GPT-4o (choices)}
  \end{subfigure}
  \vspace{0.4em}
  \begin{subfigure}[t]{0.48\linewidth}
    \includegraphics[width=\linewidth]{fig2_confusion_deepseek-r1.png}
    \caption{DeepSeek-R1 (baseline)}
  \end{subfigure}
  \hfill
  \begin{subfigure}[t]{0.48\linewidth}
    \includegraphics[width=\linewidth]{fig2_confusion_llama-3.3-70b.png}
    \caption{Llama 3.3 70B (baseline)}
  \end{subfigure}
  \caption{Row-normalized confusion matrices for all models at their best prompt.}
  \label{fig:confusion}
\end{figure}

\paragraph{Prompt effects.}
CoT degrades performance for all four models, significantly so for Claude
($-13.2$pp, $p_{\text{holm}} < 0.001$), GPT-4o
($-50.6$pp, $p_{\text{holm}} < 0.001$), and Llama
($-12.2$pp, $p_{\text{holm}} < 0.001$); the baseline vs.\ CoT gap
for DeepSeek does not survive correction ($p_{\text{holm}} = 0.143$).
The choices prompt helps Claude ($+4.3$pp, $p_{\text{holm}} < 0.001$)
but hurts DeepSeek ($-19.0$pp, $p_{\text{holm}} < 0.001$)
and has no significant effect on GPT-4o ($p_{\text{holm}} = 0.685$).
Odds ratios and contingency counts for all comparisons are in Appendix~\ref{app:stats}.

\paragraph{ROT13 memorization.}
Figure~\ref{fig:memorization} plots ROT13 accuracy against Caesar shift variants.
Claude and GPT-4o score 100\% on ROT13 while dropping substantially on
Caesar-7 (28.6\% and 81.6\%) and Caesar-19 (82.7\% and 22.4\%),
gaps that are statistically significant ($p < 0.001$, chi-squared).
Llama shows the most extreme effect: 100\% on ROT13 but 2.0\% on Caesar-7
and 0.0\% on Caesar-19.
DeepSeek shows the opposite pattern---38.8\% on ROT13 but 94.9\% on Caesar-3---
suggesting its training data distribution differs substantially.

\begin{figure}[H]
  \centering
  \includegraphics[width=0.85\linewidth]{fig3_memorization.png}
  \caption{ROT13 vs.\ Caesar shift accuracy by model. For most models, ROT13
           accuracy is inflated relative to other shifts, indicating memorization.}
  \label{fig:memorization}
\end{figure}

\paragraph{Plaintext length.}
Accuracy is largely flat across short ($\leq$62 chars), medium (63--67),
and long ($>$67) plaintexts (Figure~\ref{fig:length}), suggesting length
is not a meaningful factor within our tested range.

\begin{figure}[H]
  \centering
  \includegraphics[width=0.85\linewidth]{fig4_by_length.png}
  \caption{Accuracy by plaintext length tercile. Performance is largely invariant to length.}
  \label{fig:length}
\end{figure}

% ============================================================
\section{Limitations}

\paragraph{ROT13 memorization evidence.}
The elevated ROT13 accuracy relative to other Caesar shifts is \emph{consistent}
with name-level memorization but does not rule out alternative explanations,
such as ROT13 appearing more frequently in pretraining corpora independent of
the cipher label.
Definitive evidence would require controls that are beyond the scope of this
study: (1)~additional Caesar shifts (e.g., 11, 17, 23) to rule out
shift-specific difficulty; (2)~shuffled-alphabet substitution ciphers presented
under the ROT13 label to isolate label recognition from structural recognition;
(3)~adversarial encodings that resemble ROT13 superficially but differ
structurally; and (4)~held-out synthetic cipher names absent from pretraining data.
We treat the memorization finding as a hypothesis supported by converging evidence
rather than a confirmed causal claim.

\paragraph{Benchmark scope.}
All plaintexts are English sentences of moderate length (43--80 characters).
Generalization to other languages, longer texts, or structured data
(e.g., code, JSON) is untested.
The parser resolves common aliases but may misclassify edge-case responses;
the parse rate and raw responses are released (Appendix~\ref{app:data}) for
independent reanalysis.

% ============================================================
\section{Discussion}

CoT's consistent negative effect suggests cipher identification relies on
holistic pattern matching rather than step-by-step deduction---forcing
explicit reasoning disrupts this process.
The divergent effect of the choices prompt across models indicates that
constrained answer spaces interact non-trivially with model architecture
and instruction-following behavior.
The near-perfect encoding scores (Base64, Morse, Bacon) reflect character-set
recognition rather than genuine cryptanalysis; benchmarks should be cautious
about inflating reported numbers with such trivially-distinguished encodings.
Llama's Vigen\`{e}re bias and the universal ROT13 inflation suggest that
LLM cipher identification is driven more by training-data frequency than
by reasoning from cryptographic principles.

% ============================================================
\section{Conclusion}

LLMs can identify classical ciphers well above chance zero-shot, but their
performance reflects memorization and surface-level pattern matching more
than cryptanalytic reasoning. Chain-of-thought prompting reliably hurts;
prompt effects are model-dependent; and encoding schemes are trivially
identified. These findings have implications for both security applications
and the study of emergent capabilities in LLMs.

% ============================================================
\appendix

\section{Prompt Templates}
\label{app:prompts}

All models are queried with no system or developer prompt.
Three prompt variants are used (variable \texttt{\{ciphertext\}} is substituted at inference time):

\paragraph{Baseline.}
\begin{quote}\small
\texttt{What encryption method was used to produce the following ciphertext?}\\
\texttt{\{ciphertext\}}\\
\texttt{Answer with just the cipher name.}
\end{quote}

\paragraph{Choices.}
\begin{quote}\small
\texttt{What encryption method was used to produce the following ciphertext?}\\
\texttt{\{ciphertext\}}\\
\texttt{Choose one: Caesar, Atbash, Vigenere, ROT13, Rail Fence, Reverse, Base64, Morse Code, Bacon Cipher}\\
\texttt{Answer with just the cipher name.}
\end{quote}

\paragraph{Chain-of-thought (CoT).}
\begin{quote}\small
\texttt{What encryption method was used to produce the following ciphertext?}
\texttt{Think step by step about the patterns you observe, then give your final answer.}\\
\texttt{\{ciphertext\}}\\
\texttt{Analysis:}
\end{quote}

\section{Model and Inference Details}
\label{app:models}

\begin{table}[H]
\centering
\small
\begin{tabular}{llll}
\toprule
\textbf{Name} & \textbf{Provider} & \textbf{Model ID} & \textbf{Endpoint} \\
\midrule
Claude Sonnet 4.6 & Anthropic  & \texttt{claude-sonnet-4-6}                          & Messages API \\
GPT-4o            & OpenAI     & \texttt{gpt-4o}                                     & Chat Completions API \\
DeepSeek-R1       & DeepSeek   & \texttt{deepseek-reasoner}                          & Chat API \\
Llama 3.3 70B     & Together AI & \texttt{meta-llama/Llama-3.3-70B-Instruct-Turbo}   & Inference API \\
\bottomrule
\end{tabular}
\caption{Model identifiers and providers. All models queried at \texttt{temperature=0}, \texttt{max\_tokens=500}, no system prompt. Experiments conducted on April 25, 2025.}
\end{table}

\section{Answer Parser}
\label{app:parser}

Responses are normalized to lowercase and parsed by a rule-based extractor
in the following priority order:
(1) search for explicit answer markers (\texttt{final answer:}, \texttt{the cipher is}, etc.)
working from the end of the response;
(2) check the last three non-empty lines;
(3) check the first line;
(4) scan the full response for any valid cipher name or known alias.
A fixed alias table maps common variants (e.g., \texttt{rot-13}$\to$\texttt{rot13},
\texttt{vigen\`{e}re}$\to$\texttt{vigenere}) to canonical labels.
Responses that cannot be mapped are scored as incorrect.
The full alias table and parser source are released with the dataset.

\section{Data and Reproducibility}
\label{app:data}

The benchmark dataset (\texttt{data/dataset.jsonl}), raw model predictions
(\texttt{results/*.jsonl}), all figures, and all statistical output are
released at \url{https://github.com/hynoph/cipher-identification}.
Each prediction record contains: sample ID, true cipher label,
predicted cipher label, raw model response, parseability flag, and correctness flag.

\section{Statistical Test Details}
\label{app:stats}

McNemar's test is applied to matched prediction pairs (same sample, same prompt).
The exact binomial variant is used when the number of discordant pairs is below 25;
otherwise the asymptotic $\chi^2$ variant with continuity correction is used.
All pairwise McNemar tests (18 model comparisons $\times$ 3 prompts + 12 prompt
comparisons) are corrected jointly using the Holm--Bonferroni procedure.
Effect sizes for McNemar tests are reported as odds ratios ($b_{01}/b_{10}$,
where $b_{01}$ = cases where B is correct and A is not).
Effect sizes for $\chi^2$ tests are reported as $\phi = \sqrt{\chi^2 / n}$.
Contingency counts ($b_{01}$, $b_{10}$) for all McNemar tests are provided in
\texttt{results/stats\_mcnemar\_models.csv} and \texttt{results/stats\_mcnemar\_prompts.csv}.

% ============================================================
\bibliographystyle{plainnat}
\begin{thebibliography}{99}

\bibitem[Bubeck et al.(2023)]{bubeck2023sparks}
Bubeck, S., et al. (2023).
Sparks of artificial general intelligence: Early experiments with GPT-4.
\textit{arXiv:2303.12712}.

\bibitem[Carlini et al.(2021)]{carlini2021extracting}
Carlini, N., et al. (2021).
Extracting training data from large language models.
\textit{USENIX Security}.

\bibitem[Kojima et al.(2022)]{kojima2022large}
Kojima, T., et al. (2022).
Large language models are zero-shot reasoners.
\textit{NeurIPS}.

\bibitem[McNemar(1947)]{mcnemar1947note}
McNemar, Q. (1947).
Note on the sampling error of the difference between correlated proportions.
\textit{Psychometrika}, 12(2), 153--157.

\bibitem[Sinkov(1966)]{sinkov1966elementary}
Sinkov, A. (1966).
\textit{Elementary Cryptanalysis}.
Mathematical Association of America.

\bibitem[Wei et al.(2022)]{wei2022chain}
Wei, J., et al. (2022).
Chain-of-thought prompting elicits reasoning in large language models.
\textit{NeurIPS}.

\end{thebibliography}

\end{document}
